% Worked Examples: Concept 3.8.1 - Sensitivity Functions
\documentclass[11pt,a4paper]{article}
\usepackage{worked-examples-template}

\title{\textbf{Worked Examples}\\[0.5em]
\large Concept 3.8.1: Sensitivity Functions\\[0.3em]
\normalsize \coursesubtitle{} -- \coursetitle}
\author{Assoc.\ Prof.\ William Robertson \and Dr Sean McGowan}
\date{\today}

\begin{document}

\maketitle
\tableofcontents
\newpage

\section{Concept 3.8.1: Sensitivity Functions}

\subsection{Example 1: Computing the Gang of Four Transfer Functions}

Understanding the Gang of Four transfer functions is fundamental to analyzing feedback system performance and robustness.

\paragraph{Problem:} 
Consider a basic feedback control system with process $P(s) = \frac{5}{s(s+2)}$ and controller $C(s) = 3(1 + \frac{1}{s})$. The feedforward gain is $F(s) = 1$.

(a) Calculate the loop transfer function $L(s) = P(s)C(s)$.

(b) Determine the four fundamental sensitivity functions:
\begin{itemize}
    \item Sensitivity function: $S = \frac{1}{1+PC}$
    \item Complementary sensitivity function: $T = \frac{PC}{1+PC}$
    \item Load sensitivity function: $PS = \frac{P}{1+PC}$
    \item Noise sensitivity function: $CS = \frac{C}{1+PC}$
\end{itemize}

(c) Verify that $S + T = 1$.

(d) Evaluate each function at DC ($s = 0$) and interpret the results for disturbance rejection.

\paragraph{Solution:}

\textbf{Step 1: Calculate the loop transfer function}

The controller is a PI controller:
\[
C(s) = 3\left(1 + \frac{1}{s}\right) = 3\cdot\frac{s+1}{s} = \frac{3(s+1)}{s}
\]

The loop transfer function is:
\[
L(s) = P(s)C(s) = \frac{5}{s(s+2)} \cdot \frac{3(s+1)}{s} = \frac{15(s+1)}{s^2(s+2)}
\]

\textbf{Step 2: Calculate the sensitivity function $S$}

The sensitivity function is:
\[
S(s) = \frac{1}{1 + L(s)} = \frac{1}{1 + \frac{15(s+1)}{s^2(s+2)}}
\]

Finding a common denominator:
\[
S(s) = \frac{s^2(s+2)}{s^2(s+2) + 15(s+1)} = \frac{s^3 + 2s^2}{s^3 + 2s^2 + 15s + 15}
\]

Factoring the numerator and denominator:
\[
S(s) = \frac{s^2(s+2)}{s^3 + 2s^2 + 15s + 15}
\]

\textbf{Step 3: Calculate the complementary sensitivity function $T$}

The complementary sensitivity function is:
\[
T(s) = \frac{L(s)}{1 + L(s)} = \frac{15(s+1)}{s^3 + 2s^2 + 15s + 15}
\]

Alternatively, we can use $T = 1 - S$:
\[
T(s) = 1 - S(s) = \frac{s^3 + 2s^2 + 15s + 15 - s^3 - 2s^2}{s^3 + 2s^2 + 15s + 15} = \frac{15s + 15}{s^3 + 2s^2 + 15s + 15} = \frac{15(s+1)}{s^3 + 2s^2 + 15s + 15}
\]

\textbf{Step 4: Calculate the load sensitivity function $PS$}

The load sensitivity function is:
\[
PS(s) = P(s)S(s) = \frac{5}{s(s+2)} \cdot \frac{s^2(s+2)}{s^3 + 2s^2 + 15s + 15}
\]

Simplifying:
\[
PS(s) = \frac{5s}{s^3 + 2s^2 + 15s + 15}
\]

\textbf{Step 5: Calculate the noise sensitivity function $CS$}

The noise sensitivity function is:
\[
CS(s) = C(s)S(s) = \frac{3(s+1)}{s} \cdot \frac{s^2(s+2)}{s^3 + 2s^2 + 15s + 15}
\]

Simplifying:
\[
CS(s) = \frac{3s(s+1)(s+2)}{s^3 + 2s^2 + 15s + 15} = \frac{3s(s^2+3s+2)}{s^3 + 2s^2 + 15s + 15} = \frac{3s^3 + 9s^2 + 6s}{s^3 + 2s^2 + 15s + 15}
\]

\textbf{Step 6: Verify $S + T = 1$}

\[
S(s) + T(s) = \frac{s^3 + 2s^2}{s^3 + 2s^2 + 15s + 15} + \frac{15s + 15}{s^3 + 2s^2 + 15s + 15}
\]
\[
= \frac{s^3 + 2s^2 + 15s + 15}{s^3 + 2s^2 + 15s + 15} = 1 \quad \checkmark
\]

\textbf{Step 7: Evaluate at DC and interpret}

At $s = 0$ (DC or steady-state):

\textbf{Sensitivity function:}
\[
S(0) = \frac{0}{15} = 0
\]
This means load disturbances are \emph{completely attenuated} at steady state. This is a key benefit of integral action.

\textbf{Complementary sensitivity function:}
\[
T(0) = \frac{15}{15} = 1
\]
This confirms $S(0) + T(0) = 0 + 1 = 1$ and indicates perfect tracking of DC reference signals.

\textbf{Load sensitivity function:}
\[
PS(0) = \frac{0}{15} = 0
\]
Load disturbances have zero effect on the output at steady state.

\textbf{Noise sensitivity function:}
\[
CS(0) = \frac{0}{15} = 0
\]
Measurement noise at DC does not affect the control signal.

\textbf{Key insights:}
\begin{itemize}
\item The Gang of Four provides a complete characterization of system performance
\item The constraint $S + T = 1$ creates fundamental trade-offs: we cannot make both small simultaneously
\item Integral action (pole at $s=0$ in controller) ensures $S(0) = 0$ for perfect steady-state disturbance rejection
\item At low frequencies, we want $|S|$ small (good disturbance rejection) and $|T|$ close to 1 (good tracking)
\item At high frequencies, we want $|T|$ small (noise attenuation) which means $|S|$ approaches 1
\end{itemize}

\subsection{Example 2: Analyzing the Sensitivity and Complementary Sensitivity Trade-off}

The relationship $S + T = 1$ creates fundamental limitations on achievable performance.

\paragraph{Problem:}
Consider a feedback system with process $P(s) = \frac{1}{s+1}$ and proportional controller $C(s) = k$.

(a) Derive expressions for $S(s)$ and $T(s)$ in terms of the gain $k$.

(b) Calculate the maximum values of $|S(\ii\omega)|$ and $|T(\ii\omega)|$ over all frequencies for $k = 2$, $k = 5$, and $k = 10$.

(c) Plot the magnitude Bode plots of $S$ and $T$ for $k = 5$ and identify the crossover frequency where $|S| = |T|$.

(d) Explain the fundamental trade-off between disturbance rejection and noise amplification.

\paragraph{Solution:}

\textbf{Step 1: Derive $S(s)$ and $T(s)$}

The loop transfer function is:
\[
L(s) = P(s)C(s) = \frac{k}{s+1}
\]

The sensitivity function is:
\[
S(s) = \frac{1}{1 + L(s)} = \frac{1}{1 + \frac{k}{s+1}} = \frac{s+1}{s+1+k} = \frac{s+1}{s+(1+k)}
\]

The complementary sensitivity function is:
\[
T(s) = \frac{L(s)}{1 + L(s)} = \frac{k/(s+1)}{1 + k/(s+1)} = \frac{k}{s+1+k} = \frac{k}{s+(1+k)}
\]

We can verify: $S(s) + T(s) = \frac{s+1}{s+(1+k)} + \frac{k}{s+(1+k)} = \frac{s+1+k}{s+(1+k)} = 1$ ✓

\textbf{Step 2: Frequency response analysis}

At frequency $\omega$, substitute $s = \ii\omega$:
\[
S(\ii\omega) = \frac{\ii\omega + 1}{\ii\omega + (1+k)}
\]
\[
T(\ii\omega) = \frac{k}{\ii\omega + (1+k)}
\]

The magnitudes are:
\[
|S(\ii\omega)| = \frac{\sqrt{\omega^2 + 1}}{\sqrt{\omega^2 + (1+k)^2}}
\]
\[
|T(\ii\omega)| = \frac{k}{\sqrt{\omega^2 + (1+k)^2}}
\]

\textbf{Step 3: Maximum values for different gains}

\textbf{For $k = 2$:} $(1+k) = 3$

At low frequencies ($\omega \to 0$):
\[
|S(\ii \cdot 0)| = \frac{1}{3} = 0.333, \quad |T(\ii \cdot 0)| = \frac{2}{3} = 0.667
\]

At high frequencies ($\omega \to \infty$):
\[
|S(\ii\omega)| \to \frac{\omega}{\omega} = 1, \quad |T(\ii\omega)| \to \frac{k}{\omega} \to 0
\]

Since $|S|$ is monotonically increasing from $1/3$ to $1$: $M_s = \max |S| = 1$

Since $|T|$ is monotonically decreasing from $2/3$ to $0$: $M_t = \max |T| = 2/3 = 0.667$

\textbf{For $k = 5$:} $(1+k) = 6$
\[
|S(0)| = \frac{1}{6} = 0.167, \quad M_s = 1
\]
\[
|T(0)| = \frac{5}{6} = 0.833, \quad M_t = 0.833
\]

\textbf{For $k = 10$:} $(1+k) = 11$
\[
|S(0)| = \frac{1}{11} = 0.091, \quad M_s = 1
\]
\[
|T(0)| = \frac{10}{11} = 0.909, \quad M_t = 0.909
\]

\textbf{Step 4: Crossover frequency for $k = 5$}

The crossover frequency $\omega_c$ is where $|S(\ii\omega_c)| = |T(\ii\omega_c)|$.

Since $S + T = 1$ and both are real at the crossover:
\[
|S(\ii\omega_c)| = |T(\ii\omega_c)| = \frac{1}{\sqrt{2}} \approx 0.707 \text{ (or } -3 \text{ dB)}
\]

Setting $|S(\ii\omega_c)| = 1/\sqrt{2}$:
\[
\frac{\sqrt{\omega_c^2 + 1}}{\sqrt{\omega_c^2 + 36}} = \frac{1}{\sqrt{2}}
\]

Squaring both sides:
\[
\frac{\omega_c^2 + 1}{\omega_c^2 + 36} = \frac{1}{2}
\]

Cross-multiplying:
\[
2(\omega_c^2 + 1) = \omega_c^2 + 36
\]
\[
2\omega_c^2 + 2 = \omega_c^2 + 36
\]
\[
\omega_c^2 = 34 \quad \Rightarrow \quad \omega_c = \sqrt{34} \approx 5.83 \text{ rad/s}
\]

\textbf{Step 5: Interpretation of trade-offs}

\textbf{At low frequencies ($\omega \ll 1+k$):}
\begin{itemize}
\item $|S| \approx \frac{1}{1+k}$ is small: good disturbance rejection
\item $|T| \approx \frac{k}{1+k}$ is close to 1: good reference tracking
\item Higher gain $k$ improves both disturbance rejection and tracking
\end{itemize}

\textbf{At high frequencies ($\omega \gg 1+k$):}
\begin{itemize}
\item $|S| \to 1$: no disturbance attenuation (but high-frequency disturbances are rare)
\item $|T| \to 0$: good noise attenuation (important since noise is typically high-frequency)
\item The roll-off of $|T|$ prevents measurement noise from affecting the control signal
\end{itemize}

\textbf{At the crossover frequency $\omega_c$:}
\begin{itemize}
\item This is the transition frequency between disturbance-dominated and noise-dominated regions
\item Bandwidth is approximately at this frequency
\item Higher gain $k$ pushes $\omega_c$ higher, increasing bandwidth but making the system more sensitive to high-frequency effects
\end{itemize}

\textbf{Key insights:}
\begin{itemize}
\item The constraint $S + T = 1$ means we cannot simultaneously achieve $|S| < 1$ and $|T| < 1$ at the same frequency
\item At each frequency, we must choose between disturbance rejection (small $|S|$) and noise attenuation (small $|T|$)
\item Fortunately, disturbances are typically low-frequency and noise is typically high-frequency, so we can make $|S|$ small at low $\omega$ and $|T|$ small at high $\omega$
\item The crossover frequency represents the bandwidth of the control system
\item Increasing controller gain improves low-frequency performance but increases bandwidth, making the system more responsive to high-frequency noise
\end{itemize}
\end{document}
